\documentclass[11pt]{article}
\usepackage[T1]{fontenc}
\usepackage{amsmath,amssymb}
\usepackage[margin=32mm]{geometry}
\title{Local replacement for Lemma 3.20, cross-degree 2--4}
\date{9 October 2026}
\begin{document}
\maketitle
\noindent
\textbf{Proposed substitution for items (3) and (4) only.}
For \(w\in V'=V(H)\setminus V(G)\), let
\(S_w=N_H(w)\cap V(G)\). Every automorphism of \(G=R_3\)
extends to \(H\) preserving the two halves, so the
multiset of nine sets \(S_w\) is invariant under
\(A=\operatorname{Aut}(R_3)\).
The orbit sizes of \(A\) on 2-, 3-, and 4-subsets are,
respectively,
\[
(18,18),\qquad (6,6,36,36),\qquad
(9,9,36,36,36).
\]
As \(|S_w|=d_{\rm cross}\), the first two lists exclude
\(d_{\rm cross}\in\{2,3\}\). If \(d_{\rm cross}=4\),
the multiset is one of the two orbits of size nine,
consisting of the neighbourhoods or the non-neighbourhoods
of vertices of \(R_3\). We thus obtain an equivariant
identification of \(V'\) with \(V(G)\). Under this
identification, \(H[V']\) is \(R_3\) or its complement.
In the latter case the numbers of triangles through
vertices of the two halves are 10 and 12, contradicting
transitivity. The former case yields exactly the two
graphs in (3) and (4) of the classification, according
to the chosen four-set orbit.

\medskip
\noindent
\textbf{Audit boundary.} The orbit lists and triangle
counts are independently checked by
\texttt{checks/section3\_r3\_orbits.py}.
No original manuscript TeX is modified, and no claim of
Lean verification of this passage is made.
\end{document}
